\documentclass[10pt, letterpaper, onecolumn]{article}

% --- Preamble: Packages and Settings ---
\usepackage[utf8]{inputenc} % Set input encoding to UTF-8
\usepackage[T1]{fontenc}    % Use T1 font encoding
\usepackage{amsmath}        % For advanced math environments
\usepackage{graphicx}       % For including images
\usepackage[margin=1in]{geometry} % Set margins to 1 inch

% Optional: For better spacing and line breaks in the bibliography
\usepackage{url}
\usepackage{hyperref} % For clickable links in the PDF

% Set document title, author, and date
\title{Your Article Title Here}
\author{Your Name}
\date{\today}

% --- Document Start ---
\begin{document}

% Generates the title, author, and date
\maketitle

% Abstract (Optional)
\begin{abstract}
This is the abstract of your article. It provides a brief summary of the content, context, and results described in the document.

OK, what if I start adding stuff to the copy outside of emacs.
Does that kinda work?
\end{abstract}

% --- Main Body ---

\section{Introduction}
Start your article here. Use the $\backslash$section, $\backslash$subsection, and $\backslash$subsubsection commands to organize your content.

\subsection{Background}
Letter paper is a standard size for documents in the U.S. and Canada, measuring $8.5 \times 11$ inches. By using the $\backslash$usepackage[margin=1in]{geometry} package, we ensure the text fits nicely on the page.

\section{Methods}
This section describes the methodology. Here is an example of an equation:
$$
E = mc^2
$$
You can use inline math like $\sqrt{x^2+y^2}$ within your text.

\section{Results and Discussion}
To include an image, place the file (e.g., `figure.png`) in the same folder as your `.tex` file and use the `figure` environment.

%\begin{figure}[h]
%    \centering
%    \includegraphics[width=0.8\textwidth]{figure.png}
%    \caption{A descriptive caption for your figure.}
%    \label{fig:myfig}
%\end{figure}

For references, you can use the $\backslash$cite command and a bibliography style. For instance, see a general reference \cite{Knuth:1984}.

\section*{Conclusion}
The conclusion summarizes the main points and often suggests future work. Use $\backslash$section*{} for an unnumbered section title.

% --- Bibliography ---

\begin{thebibliography}{9}

\bibitem{Knuth:1984}
Donald E. Knuth.
\emph{The \TeX{}book}.
Addison-Wesley Professional, 1984.

\end{thebibliography}

\end{document}
